\documentclass[a4paper,french, 12pt]{article}
\usepackage[margin=1in]{geometry} 
\usepackage{amsmath}
\usepackage[many]{tcolorbox}
\usetikzlibrary{calc}
\usepackage{amssymb}
\usepackage{amsthm}
\usepackage{lastpage}
\usepackage{fancyhdr}
\usepackage{fancybox}
\usepackage{dsfont}
\usepackage{accents}
\usepackage{xcolor}
\usepackage{tikz}
\usetikzlibrary{arrows,calc,fit}
\tikzset{box/.style={draw, rectangle, thick, node distance=6em, text width=6em, text centered, minimum height=3.5em , minimum width = 8em , fill=blue!10}}
\tikzset{container/.style={draw, rectangle, dashed, inner sep=1em}}
\tikzset{line/.style={draw, thick, -latex'}}
\pagestyle{fancy}
\setlength{\headheight}{40pt}
\usepackage{stmaryrd}
\usepackage[french]{babel}
\usepackage{enumitem}
\usepackage{listings}
\geometry{hmargin=2.5cm,vmargin=1.5cm}
\renewcommand*{\overrightarrow}[1]{\vbox{\halign{##\cr 
			\tiny\rightarrowfill\cr\noalign{\nointerlineskip\vskip1pt} 
			$#1\mskip2mu$\cr}}}
\DeclareMathOperator{\Ima}{Im}
\DeclareMathOperator{\Id}{Id}
\DeclareMathOperator{\R}{\mathbb{R}}
\DeclareMathOperator{\N}{\mathbb{N}}
\DeclareMathOperator{\Z}{\mathbb{Z}}
\DeclareMathOperator{\C}{\mathbb{C}}
\DeclareMathOperator{\Q}{\mathbb{Q}}
\DeclareMathOperator{\K}{\mathbb{K}}
\let\ker\relax\DeclareMathOperator{\ker}{Ker}
\definecolor{myblue}{RGB}{65,100,250}

\tcbset{mystyle/.style={
		breakable,
		enhanced,
		outer arc=0pt,
		arc=0pt,
		colframe=myblue,
		colback=myblue!20,
		attach boxed title to top left,
		boxed title style={
			colback=myblue,
			outer arc=0pt,
			arc=0pt,
			top=3pt,
			bottom=3pt,
		},
		fonttitle=\sffamily
	}
}
\newtcolorbox[auto counter,number within=section]{Partie}[1][]{
	mystyle,
	title=Partie~\thetcbcounter,
	overlay unbroken and first={
		\path
		let
		\p1=(title.north east),
		\p2=(frame.north east)
		in
		node[anchor=west,font=\sffamily,color=myblue,text width=\x2-\x1] 
		at (title.east) {#1};
	}
}
\newtcolorbox[auto counter]{assumption}[1][]{
	mystyle,
	colback=white,
	rightrule=0pt,
	toprule=0pt,
	title=Assumption SLR.\thetcbcounter,
	overlay unbroken and first={
		\path
		let
		\p1=(title.north east),
		\p2=(frame.north east)
		in
		node[anchor=west,font=\sffamily,color=myblue,text width=\x2-\x1] 
		at (title.east) {#1};
	}
}

\setlength{\headheight}{2cm}
\setlength{\headsep}{0.6cm}
\setlength{\textheight}{26.cm} % hauteur corps texte = \textheight
%
% Largeur du texte:
\setlength{\textwidth}{17,5cm}
\setlength{\oddsidemargin}{-0.5cm}
\setlength{\evensidemargin}{0.cm}
\setlength{\marginparsep}{0.cm} 

\begin{document}
	\lhead{Thomas Arocena IM} 
	\rhead{\today} 
	\cfoot{\thepage\ / \pageref{LastPage}}
	\noindent \begin{center}
		\shadowbox{\textbf{Parcours recherche}}
		\end{center}
		
		\noindent\textbf{\underline{Simplification et choix du jeu d'instructions}}\\

		Dans un premier temps, j'utilise des hypothèses simplificatrices :
		\begin{enumerate}
		\item Les registres ont deux états possibles : initialisé ou non. Les registres initialisés contiennent une valeur entière.
		\item La pile (\textit{Stack}) ne contient que des offsets pour les appels de fonctions (\textit{function calls}).
		\item Le jeu d'instructions est réduit.
		\item Un programme passe 5 entiers en arguments stockés dans $R_1, R_2, R_3 , R_4,  R_5$ en entrée.
		\item L'état abstrait contient : \begin{enumerate}
		\item L'état du registre : initialisé ou non.
		\item Un intervalle de valeurs.
		\item Un nombre tristate qui représente la connaissance que l'on a sur les bits.\\
		\end{enumerate}
		\end{enumerate}

		Le but de ces hypothèses est d'éviter l'arithmétique des pointeurs, qui implique une formalisation laborieuse de la mémoire.
		 Se limiter dans le jeu d'instructions permet d'analyser le fonctionnement de programmes simples dans un premier temps, puis d'ajouter des instructions progressivement pour tendre vers des programmes réels.
		 Le jeu d'instructions doit donc être réduit dans un premier temps sans être trop limité, auquel cas on ne représente pas un langage fonctionnel.\\

		Jeu d'instructions :
		\begin{itemize}
		\item Instructions arithmétiques
		\begin{itemize}
		\item Addition (\textit{Add})
		\item Soustraction (\textit{Sub})
		\item Multiplication (\textit{Mul})
		\item OU binaire (\textit{Or})
		\item ET binaire (\textit{And})
		\item Affectation (\textit{Mov})
		\end{itemize}
		\item Instructions de saut
		\begin{itemize}
		\item Saut simple (\textit{Ja})
		\item Saut sous condition d'égalité (\textit{Jeq})
		\end{itemize}
		\item Appel de fonction (\textit{Call})
		\item Renvoi de valeur (\textit{Exit})\\
		\end{itemize}

		Avec ce jeu d'instructions, on peut formaliser un petit langage qui est déjà intéressant à étudier.\\
		
		\noindent\textbf{\underline{Syntaxe}}\\

		\begin{itemize}
		\item On représente un programme comme une liste d'instructions où les instructions sont celles du jeu d'instructions mentionné précédemment.
		\item On représente les registres par 11 objets : $R_0, R_1, R_2, R_3 , R_4,  R_5,  R_6,  R_7,  R_8,  R_9,  R_{10}$\
		\item On définit un argument comme un registre ou un entier.
		\item On définit la pile comme une liste d'entiers.
		\end{itemize}
		
		\noindent\textbf{\underline{Sémantique abstraite}}\
		\begin{enumerate}
		\item \textbf{\underline{État du système :}}\\
		Un état est représenté par trois informations :
		\begin{itemize}
		\item Le compteur du programme (\textit{Program counter}), noté $pc$, représente où l'on se situe dans le programme, autrement dit l'indice dans la liste d'instructions qui représente le programme.
		\item L'information des valeurs abstraites contenues dans chaque registre.
		\item La pile.
		\end{itemize}
		Dans la suite, on notera un état du système comme suit :
		
		où $reg$ est une fonction de $Register$ vers l'ensemble des valeurs abstraites noté $AbsRegisterValue$.\\

		\item \textbf{\underline{Valeur abstraite des registres (AbsRegisterValue) :}}\\
		À chaque registre, on associe un état initialisé ou non. Pour les registres initialisés, on considère :
		\begin{itemize}
		\item L'intervalle des valeurs possibles.
		\item Le nombre tristate des bits possibles.\\
		\end{itemize}

		\item \textbf{\underline{Représentation des nombres tristates:}}\\
		De manière générale, on représente un nombre tristate avec des 0 et 1 là où l'on a de l'information, et des points d'interrogation autrement.\\
		Ex : 011\textcolor{red}{?}1\textcolor{red}{??}0\\
		Mais cette représentation n'est pas efficace pour les calculs. On considère donc plutôt 2 nombres binaires : un appelé masque, qui représente la position des inconnus par des 1 et les autres par des 0, et un appelé valeur, qui représente les bits connus et met des zéros aux inconnus.\
		Ex : 011\textcolor{red}{?}1\textcolor{red}{??}0 $\to$ (00010110 (masque), 01101000 (valeur))\\
	
		\item \textbf{\underline{Instructions arithmétiques:}}\\
		Les instructions arithmétiques ne contiennent ni saut, ni appel de fonction donc le compteur du programmes passera toujours au rang suivant.
		 De la même manière la pile restera inchangée.
		Toute instruction arithmétique prenant en source un Argument correspondant à un registre non initialisé sera rejeté.
		\begin{enumerate}
			\item Addition :\\
			Pour des intervalles $i_1 = [a,b]$ et $i_2 = [c,d]$, et un entier $\lambda$:
			$$i_1 + i_2 = [a,b] + [c,d] = [a+c, b+d]$$
			$$i_1 + \lambda = [a,b] + \lambda =[a +\lambda,b + \lambda]$$
			Pour les nombres tristates :  Soit $t_1$ et $t_2$ deux nombres tristates
			$$ s_v = t_1.v + t_2.v$$
			$$ s_m = t_1.m + t_2.m$$
			$$ S = s_v + s_m$$
			$$ q = S \oplus s_v $$
			$$ h = q \lor t_1.m \lor t_2.m $$
			$$ t_1 + t_2 = (h \text{ (masque) },s_v \& \neg h \text{ (valeur) }) $$
			où $\oplus$ est le OU exclusif, $\lor$ est le OU inclusif, $\&$ est le ET et $\neg$ est la négation.\\
			Si on additionne avec une valeur immédiate $\lambda$ le calcul est identique avec $\lambda.v = \lambda$ et $\lambda.m = 0$\\
			\textit{exemple :} $([5,7],0? ) + ([1,3],?1)$\\
			pour les intervalles $[5,7] + [1,3] = [6,10]$\\
			pour les tristates $ 0? = (01,00) $ et $?1 = (10,01)$ alors :
			$$ s_v = t_1.v + t_2.v = 00 + 01 = 01 $$
			$$ s_m = t_1.m + t_2.m = 01 + 10 = 11 $$
			$$ S = s_v + s_m = 01 + 11 = 100 $$
			$$ q = S \oplus s_v = 100 \oplus 001 = 101 $$
			$$ h = q \lor t_1.m \lor t_2.m = 101 \lor 001 \lor 010 = 111 $$
			$$ t_1 + t_2 = (h ,s_v \& \neg h) = (111, 000)$$
			On obtient $([5,7],0? ) + ([1,3],? 1) = ([6,10],???)$\\

			Soit $R_m = (\text{initalisé},i_1,t_1)$ le registre de destination.\\
			\underline{Pour une source qui est un registre  $R_k = (\text{initalisé},i_2,t_2)$ on a} :
			$$ s = \{pc, reg , stack\} \to \{pc + 1, reg', stack\} $$
			où $reg'(R_n) = reg(R_n)$ si $n\neq m$ et $reg'(R_m)=(\text{initalisé},i_1+i_2,t_1+t_2)$\\

			\underline{Pour une source qui est un registre une valeur immédiate $\lambda$ on a} :
			$$ s = \{pc, reg , stack\} \to \{pc + 1, reg', stack\} $$
			où $reg'(R_n) = reg(R_n)$ si $n\neq m$ et $reg'(R_m)=(\text{initalisé},i_1+\lambda,t_1+\lambda)$\\

			\item Soustraction :\\
			Pour les intervalles:
			$$[a,b] + [c,d] = [a-d, b-c]$$\\

			\item Multiplication :\\
			Pour les intervalles:
			$$[a,b] + [c,d] = [a*c, b*d]$$\\

			Pour les tristates :\\
			\hspace*{6cm}$S_v = (0 , t_1.v*t_2.v)$\\
			\hspace*{6cm}$S_m = (0 , 0)$\\
			\hspace*{6cm}$\text{Tant que }(t_1.v \lor t_1.m \neq 0) : $\\
			\hspace*{7cm} $\text{Si } t_1.m[0] = 1 :$\\
			\hspace*{8cm} $ S_m = S_m + (t_2.v \lor t_2.m , 0) $\\
			\hspace*{7cm} $\text{Sinon si } t_2.v[0] = 1 :$\\
			\hspace*{8cm} $ S_m = S_m + (t_2.m , 0) $\\
			\hspace*{7cm} $t_1 = \text{Rshift}(t_1,1) $\\
			\hspace*{7cm} $t_2 =\text{Lshift}(t_2,1) $\\
			\hspace*{6cm}$t_1 * t_2 = S_v + S_m$\\

			où Rshift supprime le dernier bits à droite et Lshift supprime le dernier bits à gauche.\\

			\item OU binaire :\\
			Pour les nombres tristates : 
			$$ t_1 \lor t_2 = ((t_1.v \lor t_2.m) \land (t_1.m \lor t_2.v),t_1.v \lor x.v)$$\\
			Pour les intervalles :\\
			On ne fait pas ici un calcul exacte mais on utilise les tristates associés aux intervalles pour calculer plus rapidement une abstraction :
			$$ I_1 = [a,b], t_1, I_2 = [x,y], t_2 $$
			$$ t_3 = t_1 \lor t_2 $$ 
			$$ I_1 \lor I_2 = (max(a,x), t_3.v \lor t_3.m)$$

			\item ET binaire :\\
			Pour les nombres tristates :
			$$ t_1 \wedge t_2 = (t_1.v \wedge t_2.v,(t_1.v \wedge t_1.m) \lor (t_1.m \wedge t_1.v))$$
			Pour les intervalles :\\
			Idem que pour le OU
			$$ I_1 = [a,b], t_1, I_2 = [x,y], t_2 $$
			$$ t_3 = t_1 \wedge t_2 $$
			$$ I_1 \wedge I_2 = (t_3.v, min(b,y))$$

		
			\item Affectation :\\
			Dans les deux cas on affecte juste la valeur.\\

		\end{enumerate}
		\item \textbf{\underline{Instructions de saut:}}\\
		Les instructions de saut ne contiennent pas d'appel de fonction, mais jouent sur le compteur du programme. 
		La pile restera inchangée.
		\begin{enumerate}
			\item Saut simple :\\
			Pour un offset $d$ on a :
			$$ s = \{pc, reg , stack\} \to \{pc + d, reg, stack\} $$
			\item Saut sous condition d'égalité :\\
			Dans ce cas on a une complexité supplémentaire, c'est que, dans le cas abstrait, on ne peut pas toujours décider quelle branche prendre. Il faudra alors prendre les deux branches.
			\begin{center}
			\begin{tikzpicture}[auto]
		
				\node [box, fill=blue!20] (instr1) {Saut de $k$ sous condition $P$};
				\node [box, below of=instr1] (instr2) {instruction $n$};
				\node [box, below of=instr2] (instrp) {instruction $n+k-1$};
				\node [box, below of=instrp, fill=blue!20] (instrp1) {instruction $n+k$};
				\node [box, below of=instrp1, fill=blue!20] (exit) {exit};
			
				\coordinate (middle) at ($(instr2.west)!0.5!(instrp.west)$);
				%\node [box, left of=middle, node distance=10em] (archive) {Archive};
				%\node [box, left of=archive, node distance=10em] (reporting) {Reporting};
			
				\node[container, fit=(instr2) (instrp)] (non) {};
				\node at (non.north west) [above right,node distance=0 and 0] {$\neg P$};
			
				\path [line] (instr1) -- (instr2);
				%\path [line] (resources) -- (sensors);
				\path [line] (instrp) -- (instrp1);
				\node [font=\Large] at ($(instr2)!.5!(instrp)$) {\vdots};
				\node [font=\Large] at ($(instrp1)!.5!(exit)$) {\vdots};
			
				%\path [line] (instr1) |- (instrp);
				%\path [line] (archive) |- (processing);
				%\path [line] (processing)--($(processing.south)-(0,0.5)$) -| (reporting);
			
				\draw [line] (instr1.east) -- ++(1,0) node(lowerright){} |- (instrp1.east);
				\draw [line] (instr1.west) -- ++(-1,0) node(lowerleft){} |- (instr2.west);
				%\draw [line] (lowerright |- or.east) -- (or.east -| resources.south east);
			
				%\draw[line] (archive.170)--(reporting.10);
				%\draw[line] (reporting.350)--(archive.190);
			\end{tikzpicture}
		\end{center}
		On a une condition d'égalité : $R_1 = R_2$ par exemple\\
		Si $R_1$ et $R_2$ n'ont aucune valeur concretes en commun possible, i.e si les intervalles sont disjoints, alors on est sur que la condition ne sera pas vérifié.
		Sinon on passe dans vrai l'etat ou l'on a affecter l'intersection des valeurs possibles de $R_2$ et de $R_1$ à $R_1$ et on passe à faux $R_1$ inchangé.\\
		ex : $R_1 = [0,3], R_2 = [4,5] \to $ on prend la branche faux\\
		$R_1 = [4,10], R_2 = [3,5] \to $ on prend la branche faux avec $R_1 = [4,10]$ et la branche vrai avec $R_1 = [4,5]$\\

		Intersection d'intervalles :
		$$ I_1 = [a,b], I_2 = [x,y]$$
		$$ I_1 \cap I_2 = [max(a,x),min(b,y)]$$

		Intersection de tristates :
		$$ t_1 = (v_1,m_1), t_2 = (v_2, m_2) $$
		$$ t_1 \cap t_2 = ((v_1 \lor v_2)\land \neg (m_1\land m_2),m_1\land m_2)$$

		\end{enumerate}
		\item \textbf{\underline{Appel de fonction:}}\\
		Soit un offset $d$ et $p$ le compteur du programme à l'état actuel. Alors on considère le programme composé des instructions comprises entre $p+d$ et la prochaine instruction $Exit$.
		 Ce sous-programme correspond à une fonction. \`A l'instant $p$ les valeurs de $R_1,R_2,\dots,R_5$ sont les arguments passés à la fonction.\\
		\underline{Fonctionnement de la pile :} pour cela on a besoin de s'intéresser à la pile. \` chaque appel on stocke dans la pile l'indice $p$ de l'instruction $Call$.
		 On parcourt linéairement la fonction à partir de l'indice $p+d$.Quand on a parcouru la fonction jusqu'au bout, après $Exit$, on passe au rang $p+1$ et on supprime $p$ de la pile.
		\item \textbf{\underline{Stockage des états}}\\
		A chaque instruction on associe une case. On démarre avec un etat abstrait initial dans la premiere case. On suit linéairement les instructions jusqu'a un conditionnement. Lorsque on fait un conditionnement 
		\begin{center}
			\begin{tikzpicture}[auto]
		
				\node [box, fill=blue!20] (instr1) {Si $R_1 >= 5$ goto 3};
				\node [box, below of=instr1] (instr2) {$R_1 = R_1 + 5$};
				\node [box, below of=instr2] (instrp) {$R_1 = R_1*2$};
				\node [box, below of=instrp, fill=blue!20] (exit) {exit};
			
				\coordinate (middle) at ($(instr2.west)!0.5!(instrp.west)$);
				%\node [box, left of=middle, node distance=10em] (archive) {Archive};
				%\node [box, left of=archive, node distance=10em] (reporting) {Reporting};
			
				\node[container, fit=(instr2)] (non) {};
				\node at (non.north west) [above right,node distance=0 and 0] {$R_1 < 5$};
			
				\path [line] (instr1) -- (instr2);
				\path [line] (instr2) -- (instrp);
				%\path [line] (resources) -- (sensors);
				\node [font=\Large] at ($(instrp)!.5!(exit)$) {\vdots};
			
				%\path [line] (instr1) |- (instrp);
				%\path [line] (archive) |- (processing);
				%\path [line] (processing)--($(processing.south)-(0,0.5)$) -| (reporting);
			
				\draw [line] (instr1.east) -- ++(1,0) node(lowerright){} |- (instrp.east);
				\draw [line] (instr1.west) -- ++(-1,0) node(lowerleft){} |- (instr2.west);
				%\draw [line] (lowerright |- or.east) -- (or.east -| resources.south east);
			
				%\draw[line] (archive.170)--(reporting.10);
				%\draw[line] (reporting.350)--(archive.190);
			\end{tikzpicture}
		\end{center}
		Avec en entrée par exemple $R_1 = [0;10]$, on exécute ainsi :
		\begin{enumerate}[label=\textbf{I.\arabic*}]
			\item On a pas assez d'information pour savoir quelle branche prendre. On coupe est on prend la branche sans le saut.
			 On a donc $R_1=[0,4]$ passé à l'instruction 2. On stocke $R_1 = [5,10]$ et pc = 3 pour la suite.
			\item De pc = 2, $R_1=[0,4]$ on passe à pc = 3, $R_1 = [5,9]$
			\item De pc = 3, $R_1 = [5,9]$ on passe à pc = 4, $R_1 = [10,18]$
			\item On continue jusque exit. Puis on retourne à la valeur stockés.
			\item De pc = 3, $R_1 = [5,10]$ on passe à pc = 4, $R_1 = [10,20]$. Il est important de remarqué que ici, l'instruction 3 à déja ete visité.
			 Lors du passage précedent, on avait $R_1 = [5,9]$. Si l'on avait eu par exemple $R_1 = [5,11]$ alors on aurait pu s'arreter là car on aurait déja parcouru toute les branches avec un état abstrait incluant $R_1 = [5,10]$
			\item On continue jusque exit. 
		\end{enumerate}
	\item \textbf{\underline{Pruning, concaténation des états}}
		Il suffit de vérifié que lorsqu'on arrive dans un état déja visité, l'on n'a pas déja parcouru la suite avec un état abstrait "plus grand" au sens de l'inclusion.
		Inclusion pour les intervalles :\\
		$$ I_1 = [a,b], I_2 = [x,y]$$
		$$ I_2 \subseteq I_1\Leftrightarrow (a\leq x) \wedge (y\leq b)$$  
		Inclusion pour les tristates :\\
		$$ t_1 = (v_1,m_1), t_2 = (v_2, t_2) $$
		$$ t_2 \subseteq t_1\Leftrightarrow (m_1 \cup m_2 = m_1)\wedge (v_1 = (v_2 \oplus m_1))$$
		En effet il faut que tout les bits connues de $t_1$ soit les mêmes pour $t_2$. La premiere condition assure que l'on a pas d'inconnus chez $t_2$ qui sont connu chez $t_1$. La deuxieme condition assure que les valeurs connues de $t_1$ sont les mêmes chez $t_2$\\

		On peut alors faire de la concaténation d'état:
		$$A_2 = (I_2, t_2) \subseteq A_1 = (I_1, t_1) \Leftrightarrow (I_2 \subseteq I_1)\wedge(t_2 \subseteq t_1)$$
		Dans ce cas on a $$s_2 = \{pc=k, R_1 = (I_2, t_2)\} \subseteq s_1 = \{pc=m, R_1 = (I_1, t_1)\}\Leftrightarrow (k=m)\wedge((I_2,t_2)\subseteq(I_1,t_2))$$

	\item 
	\end{enumerate}
\end{document}